% --- 题目 ---
\problem[P170-2 (25-2)]{设 3 阶矩阵 $\displaystyle \bm{A}, \bm{B}$ 满足 $\displaystyle r(\bm{AB}) = r(\bm{BA}) + 1$, 则~(~\quad~)}
\begin{tasks}(1)
  \task 方程组 $\displaystyle (\bm{A} + \bm{B})\bm{x} = \bm{0}$ 只有零解.
  \task 方程组 $\displaystyle \bm{A}\bm{x} = \bm{0}$ 与方程组 $\displaystyle \bm{B}\bm{x} = \bm{0}$ 均只有零解.
  \task 方程组 $\displaystyle \bm{A}\bm{x} = \bm{0}$ 与方程组 $\displaystyle \bm{B}\bm{x} = \bm{0}$ 没有公共非零解.
  \task 方程组 $\displaystyle \bm{ABA}\bm{x} = \bm{0}$ 与方程组 $\displaystyle \bm{BAB}\bm{x} = \bm{0}$ 有公共非零解.
\end{tasks}
\ansat{336；【十年真题】 - 考点：线性方程组的解的结构 - 2}

% --- 题目 ---
\problem[P193-8 (13-1,2,3)]{设二次型
$$f(x_1,x_2,x_3)=2 \left(a_1x_1+a_2x_2+a_3x_3\right)^2 + \left(b_1x_1+b_2x_2+b_3x_3\right)^2,$$
记$\displaystyle \bm{\alpha}=\begin{pmatrix} a_1 \\ a_2 \\ a_3 \end{pmatrix}$, $\displaystyle \bm{\beta}=\begin{pmatrix} b_1 \\ b_2 \\ b_3 \end{pmatrix}$.
\begin{enumerate}
\item[(1)] 证明二次型$\displaystyle f$对应的矩阵为$\displaystyle 2\bm{\alpha}\bm{\alpha}^{\mathrm{T}}+\bm{\beta}\bm{\beta}^{\mathrm{T}}$;
\item[(2)] 若$\displaystyle \bm{\alpha}, \bm{\beta}$正交且均为单位向量, 证明二次型$\displaystyle f$在正交变换下的标准形为$\displaystyle 2y_1^2+y_2^2$.
\end{enumerate}}
\ansat{352；【真题精选】 - 考点一：化二次型为标准型 - 8}
